% TOP_PROBLEM: {"id": 3131, "title": "Are almost all graphs determined by their spectrum?", "category_id": 3, "category": {"id": 3, "name": "graph_theory", "display_name": "Graph Theory", "description": "Problems involving graphs, networks, and their properties.", "slug": "graph-theory", "order_index": 3, "created_at": "2026-07-31T15:26:25.671Z"}, "status": "open", "problem_number": "OPG-48368", "set_id": 12}
% FIRST_POSTED: 2026-10-01
% ENTRY_KIND: statement_only
% UnsolvedMath ID 3131; source catalogue code OPG-48368
% !TeX program = lualatex
% Definition and sources only; this file is not an LLM solution attempt.
\documentclass[11pt,a4paper]{article}
\usepackage[margin=25mm]{geometry}
\usepackage{fontspec}
\setmainfont{lmroman10-regular.otf}[BoldFont=lmroman10-bold.otf,
 ItalicFont=lmroman10-italic.otf,BoldItalicFont=lmroman10-bolditalic.otf]
\usepackage{amsmath,amssymb,mathtools}
\usepackage{unicode-math}
\setmathfont{latinmodern-math.otf}
\AtBeginDocument{\let\setminus\smallsetminus}
\usepackage{xurl,hyperref}
\hypersetup{colorlinks=true,urlcolor=blue}
\setcounter{secnumdepth}{0}
\setlength{\parindent}{0pt}
\setlength{\parskip}{5pt}
\setlength{\emergencystretch}{3em}
\begin{document}
\section*{Are almost all graphs determined by their spectrum?}\label{3131}
{\small UnsolvedMath ID 3131 · OPG-48368 · Graph Theory · OpenGarden}
\subsection{Definitions and mathematical statement}
For an integer $n\geq1$, let $\mathcal G_n$ be the set of all simple
undirected graphs on the labelled vertex set $[n]=\{1,\ldots,n\}$.
The adjacency matrix $A_G$ has entry $1$ at $(u,v)$ precisely when
$uv\in E(G)$, and has zero diagonal. Its adjacency spectrum is the
multiset of its $n$ real eigenvalues, including their multiplicities.
Two graphs are cospectral when these multisets agree; equivalently,
their characteristic polynomials $\det(tI-A_G)$ agree.

A graph $G$ is determined by its adjacency spectrum if every finite
simple graph $H$ cospectral with $G$ is isomorphic to $G$. Here an
isomorphism is a bijection between the vertex sets preserving adjacency
in both directions. Relabelling a graph therefore does not count as a
different spectral realization. Equality of spectra already forces
$|V(H)|=|V(G)|$.

The question is whether
\[
 \lim_{n\to\infty}
 \frac{|\{G\in\mathcal G_n:G\text{ is determined by its spectrum}\}|}
 {2^{\binom n2}}=1.
\]
Equivalently, choose each possible edge independently with probability
$1/2$; does the probability of a nonisomorphic cospectral mate tend to
zero? This specifies the usual labelled random-graph meaning of
``almost all.'' The data supplied are only the adjacency spectrum:
neither eigenvectors nor the spectrum of the complement or Laplacian
is included in the reconstruction requirement.
\subsection{Short English statement}
Does the adjacency spectrum uniquely identify a uniformly random large graph, up to renaming its vertices, with probability tending to one?
\subsection{Sources}
\begin{itemize}
\item[S1] ULAM.AI and the UnsolvedMath Contributors, supplied problems.json, record 3131 (OPG-48368), OpenGarden collection. Catalogue identity and source transcription; CC BY 4.0.
\url{https://huggingface.co/datasets/ulamai/UnsolvedMath}.
\item[S2] Open Problem Garden, Are almost all graphs determined by their spectrum?. Original problem and discussion; the mathematical formulation below is expanded from the supplied source record.
\url{http://www.openproblemgarden.org/op/are_almost_all_graphs_determined_by_their_spectrum}.
\end{itemize}
\end{document}
